% =====================================================================
% Section 5 --- Application: the programme through the matrix
% ---------------------------------------------------------------------
% Job: multipliers, price responses, external position across cells; the
%      robust-vs-closure-sensitive table; decision heuristics for
%      transformation modellers. THIS IS THE HEADLINE THE INTRO PROMISES.
% Mindmap: N06, N05, N10.
% Status: STUB --- narrative outline v3 (2026-09-20).
% =====================================================================
\section{Application: the programme through the matrix}\label{sec:application}

% TODO: run the green-investment programme through the 5x3 matrix.
% Every table cites run_ids from matrix_5x3_v10 (matrix + 18 sectoral cells)
% and supply_etas_* (189 cells; s1 the identification vehicle).
% TABLES (linked, revision/tables/):
%   - matrix_5x3_v10_flows.md  -> the 15 matrix cells + 18 sectoral cells (Tables 1-4)
%   - supply_etas_flows.md     -> the supply arm s1 (identification vehicle)
%   - supply_etas_prog_flows.md / supply_etas_unif_flows.md -> the other two designs
% Port each to LaTeX at its point of use.

\subsection{Aggregate results}\label{subsec:aggregate}
% TODO: multipliers and price responses; the wage-regime vs elasticity
% channels (R1.7). Real GDP is the income-side index, flat by construction
% under demand-only shocks (ADR-0018); consumption is welfare. Do NOT repeat
% the old "virtually zero to ~1.5 %" GDP sentence (retired framing).

\subsection{Sectoral results}\label{subsec:sectoral}
% TODO: sectoral detail; the ladder and rigid-group cells.

\subsection{The financing wedge}\label{subsec:financing-wedge}
% TODO: bracket results (F1/F2/F3); the fixed-wage row is where financing
% bites (GAMMA/DELTA F2 vs F3: -1.53 % vs +2.26 % consumption).

\subsection{Robust vs closure-sensitive}\label{subsec:robust}
% TODO: the decision table for transformation modellers; which conclusions
% survive every closure, which flip.

% ===== CARRIED OVER FROM AN OLD DRAFT =============================
% [devised for an old draft; not in the validated submission]
% source: revised_manuscript/chapters/ch05_ces_functions.tex (commit 41144dd)
% =================================================================
% An overall inspection of the results show that in this comparison the standard Leontief approach delivers the largest GDP-estimates -- on this basis GDP is expected to grow by about 1.5\%, which is equivalent to a multiplier of 1.15 given that the initial investment impulse was specified to be 40.337 bn. €, about 1.3\% of GDP. In contrast, the Cobb-Douglas specification predicts a small reduction in GDP that results from efficiency-losses arising from the increase in green investment. These losses are due to the fact that the imposed green investment leads to a forced reallocation of intermediates and production factors towards the expanding sectors, which alters an already optimal allocation and, hence, induces inefficiencies. In turn, real output decreases slightly and the major macroeconomic impact of a green investment demand shock in this framework is to induce sector-specific inflationary pressures.
% Finally, by assessing changes in sectoral production after the demand shock in Figure \ref{fig:sectoral-changes}, we find that the output adjustment approach will lead to an expansion of almost all sectors, while the general equilibrium approach leads to a reduction of gross output in the large majority of sectors, that are not directly affected by the demand shock. Moreover, the shock also depresses gross output and final consumption associated with two of the shocked sectors (in line with Figure \ref{fig:comparison-sec4} above, which focuses on final demand only). This reproduces the finding from section \ref{subsec:CGE_elasticity} that a demand shock will depress real GDP on a sectoral level as the (comparatively lower) average expansion of shocked sectors does under no conditions fully compensate for the reduction in the production of unaffected sectors. Notwithstanding this overall pattern, the complex substitution dynamics implied the CGE approach lead to exceptional outcomes in specific sectors. The most obvious idiosyncratic outcome occurs in two sectors that are not shocked –– \textit{building construction works (sector 33 )} and \textit{civil engineering works (sector 34)} --, but, rather, massively expand their intermediate goods production at the expense of final outputs, which represents an over-proportional intersectoral crowding-out effect. In other words, intermediary goods from these sectors are in extremely high demand due to the enforced shift in final expenditures. For both of these sectors the expansion in intermediate goods production is so strong, that these sectors expand more strongly than in the Leontief case and thereby partially compensate for the lack of freely available workers in shocked sectors.

